\documentclass[aspectratio=169, 10pt]{beamer}

\input{preamble_new}

% --- Meta -------------------------------------------------------------------
\title{Practical AI for Academics}
\subtitle{Day 2 --- hands-on with your own materials}
\author{Claes B\"ackman}
\institute{Leibniz Institute for Financial Research SAFE}
\date{Day 2 --- August 28, lugano}

% ============================================================================
\begin{document}

\begin{frame}[plain]
  \titlepage
\end{frame}

%----------------------------------------------------------------------------
% AGENDA
% ----------------------------------------------------------------------------
\sourceline{}

\begin{frame}{Day 2 agenda}
\begingroup
\hyphenpenalty=10000 \exhyphenpenalty=10000
\begin{tikzpicture}[
  block/.style={rectangle, draw=mcknavy, line width=0.6pt, fill=mcknavy!5,
    minimum width=1.7cm, minimum height=1.4cm, align=center,
    font=\sffamily\small\bfseries, text=mckdeep, rounded corners=2pt},
  arr/.style={-{Stealth[length=2mm]}, mckblue, line width=1pt},
  capt/.style={font=\scriptsize\sffamily, text=mckgrey, align=center,
    text width=1.85cm, inner sep=1pt},
]
\node[block] (b1) at (0,0)     {1\\[2pt] \footnotesize Intro};
\node[block] (b2) at (1.92,0)  {2\\[2pt] \footnotesize Referee\\ report};
\node[block] (b3) at (3.84,0)  {3\\[2pt] \footnotesize Secure\\ data};
\node[block] (b4) at (5.76,0)  {4\\[2pt] \footnotesize Writing};
\node[block] (b5) at (7.68,0)  {5\\[2pt] \footnotesize Paper\\ \& code};
\node[block] (b6) at (9.6,0)   {6\\[2pt] \footnotesize Slides};
\node[block] (b7) at (11.52,0) {7\\[2pt] \footnotesize Repli-\\ cation};
\draw[arr] (b1) -- (b2);
\draw[arr] (b2) -- (b3);
\draw[arr] (b3) -- (b4);
\draw[arr] (b4) -- (b5);
\draw[arr] (b5) -- (b6);
\draw[arr] (b6) -- (b7);

\node[capt] at (0,-1.55)     {Where we left off; what today covers};
\node[capt] at (1.92,-1.55)  {Demo, then a report on your own draft};
\node[capt] at (3.84,-1.55)  {Register data; what never leaves the server};
\node[capt] at (5.76,-1.55)  {A \texttt{voice.md} from your papers; editing};
\node[capt] at (7.68,-1.55)  {Does the code do what the paper says?};
\node[capt] at (9.6,-1.55)   {Beamer decks from one style file};
\node[capt] at (11.52,-1.55) {Packaging code and data to hand over};
\end{tikzpicture}
\endgroup

\vskip0.15cm
\textbf{Slides, skills, and guides for both days:} \href{https://claesbackman.com/lugano}{\texttt{claesbackman.com/lugano}}
\end{frame}


% ============================================================================
\section{A referee report on your paper}
% ============================================================================
\sourceline{}

\begin{transitionframe}
Pick one draft you would like another pair of eyes on. Preferably a relatively close-to-finished draft. \\
\bigskip 
Then write /review-paper 
\end{transitionframe}



\begin{frame}{\texttt{/review-paper} on a real draft, eight subagents in parallel.}
\begin{columns}[t]
\column{0.48\textwidth}
\textbf{Setup}
\begin{itemize}\footnotesize
  \item One of my own papers, in its own folder
  \item \texttt{.tex} source, \texttt{.bib}, supporting materials
  \item One command: \texttt{/review-paper} with a target journal
\end{itemize}

\vskip0.15cm
\textbf{Output}
\begin{itemize}\footnotesize
  \item A structured markdown report
  \item Referee-style comments, grouped by theme
  \item Specific line and section references
\end{itemize}

\column{0.48\textwidth}
\begin{mckcardg}[mcknavy]{Eight subagents running in parallel}{demo}
\footnotesize
\begin{itemize}\setlength{\itemsep}{0pt}
  \item Spelling, grammar \& style
  \item Internal consistency \& cross-references
  \item Unsupported claims \& identification
  \item Mathematics, equations \& notation
  \item Tables, figures \& documentation
  \item Referee assessment: positioning \& fit
  \item Contribution advocate
  \item Contribution skeptic
\end{itemize}
\textcolor{mcknavy}{\itshape\scriptsize \textbf{Pattern:} eight narrow critics beat one diffuse one.}
\end{mckcardg}
\end{columns}

\end{frame}

\begin{frame}{Run a referee report on your own writing.}
\vskip0.2cm
\begin{enumerate}
  \item Open one of your own drafts in VS Code
  \item Make sure it is in markdown or LaTeX --- convert if needed
  \item Run \texttt{/review-paper-light} for a 1--2 minute version, or \texttt{/review-paper} for the full eight-agent run
  \item Read the output. It's not all correct! 
\end{enumerate}

\vskip0.4cm
\textbf{Debrief questions}
\begin{itemize}
  \item What did the report catch that you would not have seen by yourself?
  \item What did it get wrong?
  \item Which subagent gave you the most useful feedback?
\end{itemize}
\end{frame}


\begin{frame}[fragile]{A skill is a markdown file the agent loads on demand --- invoke with \texttt{/name}.}
\vskip0.1cm
\small
A skill lives at \texttt{.claude/skills/your-skill/SKILL.md} (or \texttt{\textasciitilde/.claude/skills/} for every project):
\begin{verbatim}
---
name: review-paper
description: Run an 8-agent pre-submission referee report for an academic paper
targeting a specified journal
---
You are coordinating a rigorous pre-submission review of an 
academic economics paper. You will run 8 specialized review agents 
in parallel and consolidate their findings into a structured report.

\end{verbatim}

\vskip0.2cm
Invoke with \texttt{/review-paper}. Claude loads the instructions and runs them.
\end{frame}

\begin{frame}{\texttt{CLAUDE.md} is always loaded, skills are on demand.}
\vskip0.1cm
\begin{columns}[t]
\column{0.48\textwidth}
\begin{mckcard}[mcknavy]{\texttt{CLAUDE.md} --- always loaded}
\footnotesize
\begin{itemize}
  \item Loaded into context every session
  \item Background, conventions, voice
  \item ``How I want you to work with me''
\end{itemize}
\vskip0.3em
\textcolor{mcknavy}{\itshape\scriptsize \textbf{Cost:} permanent context budget.}
\end{mckcard}

\column{0.48\textwidth}
\begin{mckcard}[mckamber]{Skill --- on demand}
\footnotesize
\begin{itemize}
  \item Loaded only when invoked with \texttt{/name}
  \item A specific workflow, with structured output
  \item ``Do this thing for me''
\end{itemize}
\vskip0.3em
\textcolor{mckamber}{\itshape\scriptsize \textbf{Cost:} zero until triggered.}
\end{mckcard}
\end{columns}

\vskip0.4cm
\begin{sowhat}
If you find yourself pasting the same prompt multiple times, you can ask Claude to make it into a skill. 
\end{sowhat}
\end{frame}




\begin{frame}{Subagents add two things the main agent cannot do alone --- parallelism and context isolation.}
\centering
\begin{tikzpicture}[scale=0.95, transform shape,
  every node/.style={font=\sffamily\small, align=center},
  box/.style={rectangle, draw=mcknavy, rounded corners=2pt, fill=mcknavy!8,
    minimum width=2.2cm, minimum height=1.0cm, inner sep=4pt, text=mckdeep},
  mainbox/.style={rectangle, draw=mcknavy, line width=0.9pt, rounded corners=2pt,
    fill=mcknavy!18, minimum width=4.5cm, minimum height=1.1cm, inner sep=5pt,
    text=mckdeep, font=\sffamily\small\bfseries},
  arrow/.style={-{Stealth[length=2mm]}, line width=0.9pt, mcknavy},
  back/.style={-{Stealth[length=2mm]}, line width=0.9pt, mckamber},
]
\node[mainbox] (m) at (0,2.6) {Main agent\\[2pt]\scriptsize\mdseries the conversation you are in};

\node[box] (a) at (-3.2,0) {\textbf{Subagent A}\\\scriptsize empirical critic};
\node[box] (b) at ( 0.0,0) {\textbf{Subagent B}\\\scriptsize theory critic};
\node[box] (c) at ( 3.2,0) {\textbf{Subagent C}\\\scriptsize methods critic};

\node[font=\scriptsize\sffamily, text=mckgrey] at (-3.2,-0.95) {\textit{fresh context}};
\node[font=\scriptsize\sffamily, text=mckgrey] at ( 0.0,-0.95) {\textit{fresh context}};
\node[font=\scriptsize\sffamily, text=mckgrey] at ( 3.2,-0.95) {\textit{fresh context}};

\draw[arrow] (m) -- (a) node[midway, left=2pt, font=\scriptsize\sffamily] {task};
\draw[arrow] (m) -- (b);
\draw[arrow] (m) -- (c);

\draw[back] (a.east) ..controls (-1.5,1.7).. (m.south west);
\draw[back] (b.north) -- (m.south);
\draw[back] (c.west) ..controls (1.5,1.7).. (m.south east)
  node[midway, right=4pt, font=\scriptsize\sffamily, text=mckamber] {summary};
\end{tikzpicture}

\vskip0.15cm
\begin{sowhat}
The \texttt{/review-paper} skill spawns \emphnavy{eight} subagents.
\end{sowhat}
\end{frame}


\begin{transitionframe}
Now launch another agent to critique the referee report you just generated --- ask Claude to write the instructions.
\end{transitionframe}


\begin{frame}{The second-opinion pattern --- spawn an adversary to check on the work. }
\textit{The main agent sees its own prior answer and has reason to defend it.}

\vskip0.2cm
\textbf{Why a fresh subagent is different}
\begin{itemize}\footnotesize
  \item No chat history: no anchoring on what you already wrote
  \item No conversational momentum: no reason to agree with the main agent
  \item Clean context: can hold the whole artifact at once
\end{itemize}

\vskip0.2cm
\textbf{When to spawn one}
\begin{itemize}\footnotesize
  \item Before accepting a long answer: \textit{``is this actually right?''}
  \item After a long refactor: \textit{``what did we break?''}
  \item When you suspect flattery: \textit{``give me the honest version''}
  \item Before sending: \textit{``read this as a hostile reader''}
  \item After doing automatic edits: \textit{``look over the changes we did''}
\end{itemize}

\vskip0.1cm
\begin{sowhat}
The phrases ``\emphnavy{spawn a subagent}'' or ``\emphnavy{launch an agent}'' are the trigger words. The main agent knows what to do and will write the instructions for you.
\end{sowhat}
\end{frame}

% \begin{frame}[fragile]{Three patterns cover most adversarial-review use cases.}
% \vskip0.2cm
% \textbf{1.\ Critique a draft as a hostile referee}
% \begin{quote}\small\itshape\color{mcknavy}
% ``Launch an agent to read \texttt{draft.tex} and give me the three strongest objections a referee would raise.''
% \end{quote}

% \textbf{2.\ Verify numbers against code}
% \begin{quote}\small\itshape\color{mcknavy}
% ``Spawn a subagent to re-derive the numbers in Table 3 from \texttt{05\_results.do} and tell me if they match the draft.''
% \end{quote}

% \textbf{3.\ Stage an internal debate}
% \begin{quote}\small\itshape\color{mcknavy}
% ``Spawn two subagents: one argues the identification strategy works, one argues it doesn't. Then summarise the disagreement.''
% \end{quote}

% \vskip0.3cm
% \begin{sowhat}
% You do not need to write the subagent's instructions. The main agent does that --- you only set the framing.
% \end{sowhat}
% \end{frame}

\begin{frame}{What the demo is meant to show}
\vskip0.2cm
\begin{enumerate}
  \item \emphnavy{Grounded output.} The report is about \emph{your} draft, not a generic checklist
  \item \emphnavy{Subagents are why this works.} Eight narrow critics catch things one diffuse reader misses
  \item \emphnavy{Not a replacement for a referee.} It is a pre-flight check before you submit
  \item \emphnavy{Transparent.} The skill is a $\sim$200-line markdown file --- fork it, modify it, make it yours
\end{enumerate}

\vskip0.4cm
\begin{sowhat}
The point is not to trust the report. The point is to make hard feedback cheap enough that you ask for it before submitting --- not after the desk reject.
\end{sowhat}
\end{frame}

\begin{frame}{Other tools generate referee reports too}
\vskip0.2cm
\small
\begin{tabular}{>{\bfseries\color{mckdeep}}lp{6.5cm}>{\itshape\color{mckgrey}}l}
\toprule
\rowcolor{mcklight}
\textbf{\textcolor{mckdeep}{Tool}} & \textbf{\textcolor{mckdeep}{What it does}} & \textbf{\textcolor{mckdeep}{Pricing}} \\
\midrule
Refine.ink                  & Referee-style report on a paper PDF      & Paid \\
coarse.ink                  & Lightweight version of the same idea     & $\approx$ free \\
Referee 2                   & Paid referee-report tool                 & Paid \\
Stanford Agentic Reviewer   & Open agentic reviewer                    & Free \\
\bottomrule
\end{tabular}


\vskip0.3cm
\begin{sowhat}
The advantage of running your own skill is that it lives next to your draft, knows your conventions, and you can edit what it checks.
\end{sowhat}
\end{frame}



%============================================================================
\section{Working with secure data}
% ============================================================================
\sourceline{}

\begin{frame}
\centering
  \emphwarn{You should never paste any sensitive data or results into Claude} \\ 
  \bigskip 
  \emphwarn{Files or chats that Claude read go to Anthropic servers}
\end{frame}


\begin{frame}{Confidentiality and data}
\begin{itemize}
  \item \textbf{Never paste} confidential microdata, register data, or unpublished co-authored work.
  \item \textbf{``Training on your data.''} Some consumer tiers may use your chats by default.
        Know the setting; prefer business/team tiers for anything sensitive.
  \item \textbf{When in doubt, abstract it.} Ask about a made-up example or use a   synthetic dataset.
  \item \textbf{Check institutional policy} before putting any work material into any AI tool.
\end{itemize}
\end{frame}

% 

\begin{frame}{Secure servers split the workflow}
\vskip0.2cm
\textbf{What Claude can still help with on your laptop}
\begin{itemize}
  \item Your \texttt{.do} files
  \item Writing and iterating on empirical strategy 
\end{itemize}
\bigskip
Claude cannot run the files or read the logs on the server, but can still be very helpful in a number of ways 

\end{frame}

\begin{frame}{Two artifacts to keep local}
\begin{columns}[t]
\column{0.48\textwidth}
\begin{mckcard}[mcknavy]{1\ \ Codebook in markdown}
\footnotesize
\begin{itemize}
  \item Convert once (\texttt{/pdf-to-markdown})
  \item Variable names, labels, value ranges, missing-value codes
  \item Agent reads it before writing \texttt{recode}, \texttt{label define}, \texttt{merge}
\end{itemize}
\end{mckcard}

\column{0.48\textwidth}
\begin{mckcard}[mcknavy]{2\ \ Simulated dataset}
\footnotesize
\begin{itemize}
  \item Same variable names, types, ranges --- not real values
  \item AER/JF/ReStud replication packages often ship with these
  \item Or generate from the codebook
\end{itemize}
\end{mckcard}
\end{columns}
\end{frame}


\begin{frame}{A workflow for secure data on servers}
\vskip0.3cm
\begin{tikzpicture}[
  wstep/.style={rectangle, draw=mcknavy, line width=0.6pt, fill=mcknavy!7,
    rounded corners=2pt, text width=3.3cm, minimum height=1.4cm, align=center,
    font=\sffamily\footnotesize, text=mckdeep},
  arr/.style={-{Stealth[length=2.5mm]}, line width=0.9pt, mckblue},
]
\node[wstep] (s1) at (0,2)    {\textbf{1.\ Set up local folder}\\\scriptsize codebook + sim data + \texttt{.do} + project description};
\node[wstep] (s2) at (4.3,2)  {\textbf{2.\ Iterate with agent}\\\scriptsize until do-file runs on sim data};
\node[wstep] (s3) at (8.6,2)  {\textbf{3.\ Copy do-file manually}\\\scriptsize run on real data};
\node[wstep] (s4) at (0,0)    {\textbf{4.\ See what worked and what you need to change}\\\scriptsize};
\node[wstep] (s5) at (4.3,0)  {\textbf{5.\ Ask for feedback or updates }\\\scriptsize ``I now want to run this regression'' };
\node[wstep] (s6) at (8.6,0)  {\textbf{6.\ Copy back }\\\scriptsize };

\draw[arr] (s1) -- (s2);
\draw[arr] (s2) -- (s3);
\draw[arr] (s3.south) -- ++(0,-0.5) -| (s4.north);
\draw[arr] (s4) -- (s5);
\draw[arr] (s5) -- (s6);
\end{tikzpicture}

\vskip0.4cm
\begin{sowhat}
The agent never sees the real data \emphwarn{(and you should never paste any result from the server into Claude)}--- but it sees everything that explains the data. That is usually enough to write correct Stata code.
\end{sowhat}
\end{frame}

% ----------------------------------------------------------------------------
% Does the simulated-data step actually work? Tested on a Danish register
% codebook. Follows directly from "Two artifacts to keep local" (artifact 2)
% and step 2 of the workflow diagram above.
% ----------------------------------------------------------------------------

% \begin{frame}{Does step 2 actually work? Tested on a Danish register codebook}
% \bigskip
% \textit{Artifact 2 above says ``or generate from the codebook''. Here is what that
% produces.}
% \bigskip

% Two codebooks (BEF population, IND income), no data, fresh session. The agent
% produced 55,000 rows per register with:
% \begin{itemize}\setlength{\itemsep}{2pt}
%   \item Correct variable names --- \texttt{PNR}, \texttt{KOEN}, \texttt{CIVST},
%         \texttt{ERHVERVSINDK\_13}, \texttt{FORM}, \texttt{QSPLINDK}
%   \item Correct value codes --- \texttt{KOEN} in \{1,2\}, \texttt{CIVST} in
%         \{U,G,F,E\}, \texttt{IE\_TYPE} in \{1,2,3\}
%   \item A working \texttt{PNR}--\texttt{AAR} merge key across both registers
% \end{itemize}
% \bigskip \pause
% \begin{sowhat}
% Code written against this runs unchanged on the real registers. The schema is
% right, so merges, reshapes, variable names and sign conventions are all tested
% before you ever enter the secure environment.
% \end{sowhat}
% \end{frame}



\begin{frame}{What synthetic data does and does not buy you}
\vskip0.1cm
\begin{columns}[t]
\column{0.49\textwidth}
\begin{mckcardg}[mckteal]{Tests that the code RUNS}{cbk}
\footnotesize
Variable names, merge keys, reshapes, type handling, sign conventions, the
shape of the panel. All the things that produce an error message. 
\end{mckcardg}

\column{0.49\textwidth}
\begin{mckcardg}[mcknavy]{Does NOT test that it is RIGHT}{cbk}
\footnotesize
Correlations, magnitudes, selection, missingness patterns. The fake data can only contain what someone thought to put in it
\end{mckcardg}
\end{columns}

\end{frame}



% ============================================================================
\section{Writing with agentic AI}
% ============================================================================
\sourceline{}


\begin{frame}{What do you actually need from AI?}
    \includegraphics[width=0.7\linewidth]{sophia_1} \\
    \pause 
    \begin{center}
          \includegraphics[width=0.7\linewidth]{sophia_2} \\

    \end{center}

\end{frame}



\begin{frame}{Most of the value is not in having AI write your text}
  I like writing, and I don't want AI to take away part of the job that I like \\ 
  % \bigskip 
  % But: I am not a native speaker, and I make mistakes \\
  % \bigskip
  % AI can be very helpful in helping me structure ideas or fix mistakes, but it is worth emphasizing that a lot of AI writing is bad. \\ 
  \bigskip
  Use it for what it is good for, and learn to ignore some of what it says 
\end{frame}

\begin{frame}{Use AI to edit your draft, don't edit the AI's draft}

\vskip0.1cm
\begin{columns}[t]
\column{0.49\textwidth}
\begin{mckcardg}[mckteal]{  Writing is thinking}{cbk}
\footnotesize
If you remove too much of the writing friction, it will be harder to figure out what your paper is about
\end{mckcardg}

\column{0.49\textwidth}
\begin{mckcardg}[mcknavy]{People are (rightfully) annoyed with AI writing}{cbk}
\footnotesize
It is now easy to spot AI writing, especially with tools like Pangram. You don't want people to discount your idea
\end{mckcardg}
\end{columns}

\end{frame}

\begin{frame}[fragile]{A \texttt{voice.md} file with your style}
\vskip0.1cm
\small
\begin{verbatim}
When editing my prose, match the voice in voice.md.
Do not use em-dashes, bullet lists, or passive
constructions unless the original does.
Do not use "it's not x, it's y"-style writing.
\end{verbatim}

\vskip0.3cm
\textbf{What goes in it}
\begin{itemize}
  \item 5--10k words of your prose --- intros, abstracts, paragraphs you are proud of
\end{itemize}

\end{frame}


\begin{frame}[fragile]{Hands-on (15 min) --- build a \texttt{voice.md} from your own writing.}
\vskip0.2cm
We use a skill written by Mihail Velikov: \texttt{voice-extractor}.

\vskip0.3cm
\begin{enumerate}
  \item Upload 4 PDFs of your own writing to a folder
  \item Call the skill: \texttt{/voice-extractor sample-papers}
  \item Read through \texttt{voice.md} and verify --- edit anything that does not sound like you
\end{enumerate}

\vskip0.4cm
\textbf{What goes in.} Intros, abstracts, paragraphs you are proud of. 5--10k words is plenty.\\
\textbf{What stays out.} Co-authored sections where the prose is not yours. AI-edited drafts.

\vskip0.3cm
\begin{sowhat}
Reference \texttt{voice.md} from \texttt{CLAUDE.md}. Every future session loads your voice into context.
\end{sowhat}
\end{frame}

\begin{frame}{\texttt{/prose-revision} fixes the order of the ideas before the choice of words.}
\vskip0.25cm
\textbf{What comes back.} A diagnosis of the idea flow first, then two or three rewrites that differ in how far they restructure, and a short note on the tradeoff between them. Claims, citations, and notation come through unchanged.

\vskip0.25cm
\textbf{Where it stops.} It will not add or drop an idea without saying so, and it is not a proofreader. Revision comes first, the typo pass comes last.

\vskip0.25cm
{\footnotesize From \url{https://x.com/alz_zyd_}}
\end{frame}


% \begin{frame}{How to tell whether your \texttt{voice.md} is working}
% \vskip0.2cm
% \begin{columns}[t]
% \column{0.48\textwidth}
% \begin{mckcard}[mckteal]{Signs it is working}
% \footnotesize
% \begin{itemize}
%   \item Edited paragraphs read more like you, less like a generic LLM
%   \item Fewer em-dashes, fewer bullet lists, fewer ``crucially'' / ``moreover''
%   \item You stop having to say ``preserve my voice'' in every prompt
%   \item Your coauthors stop noticing which paragraphs were AI-edited
% \end{itemize}
% \end{mckcard}

% \column{0.48\textwidth}
% \begin{mckcard}[mckwarn]{Signs it is not}
% \footnotesize
% \begin{itemize}
%   \item Output still sounds generic
%   \item AI tics persist after editing
%   \item You keep adding ``don't use X'' to every prompt
% \end{itemize}
% \vskip0.3em
% \textcolor{mckwarn}{\itshape\scriptsize \textbf{Fix:} sample is too short, or includes too much AI-edited prose. Add more of \emph{your} writing.}
% \end{mckcard}
% \end{columns}

% \vskip0.4cm
% \begin{sowhat}
% \texttt{voice.md} is a living file. Update it when your style shifts, or when an editor pushes you in a direction you want to keep.
% \end{sowhat}
% \end{frame}


\begin{frame}{AI detectors can tell AI generated writing}
\vskip0.1cm
The most popular current detector is \textbf{Pangram}, with very low reported false-positive rates.  \\
\bigskip 
\textbf{How they work.} A classifier trained on the fingerprint of model text. \\
\bigskip 
Does it work? This is almost irrelevant today 
\bigskip
\end{frame}



% ============================================================================
\section{Verifying the paper-code link and changes to the code}
% ============================================================================

\begin{frame}{\texttt{/review-paper-code} asks whether the code does what the paper says.}
\begin{columns}[t]
\column{0.50\textwidth}
\textbf{What it does, in four steps}
\begin{enumerate}\small\setlength{\itemsep}{1pt}
  \item Finds your main \texttt{.tex} and your code folder on its own --- Stata, R, Python
  \item Reads the paper for the sample, the specification, and the headline numbers
  \item Runs two agents at once: reproducibility, and paper-to-code mapping
  \item Reconciles the two and writes \texttt{reviews/code\_review\_report.md}
\end{enumerate}

\vskip0.15cm
\small
\textbf{Call it:} \texttt{/review-paper-code}, adding \texttt{full} to review every file it finds.

\column{0.46\textwidth}
\begin{mckcardg}[mcknavy]{What comes back}{rpc}
\footnotesize
\begin{itemize}\setlength{\itemsep}{1pt}
  \item \textbf{Reproducibility checklist:} paths, seeds, ungenerated outputs, dependencies, run order
  \item \textbf{Paper--code map:} matched, worth verifying, or not found in the files reviewed
  \item \textbf{Likely discrepancies:} where the two disagree
  \item \textbf{Next steps,} ordered by importance
\end{itemize}
\vskip0.1em
\end{mckcardg}
\end{columns}

\vskip0.1cm
\begin{sowhat}
Does Table 3 come from the regression the paper describes? Cheapest to answer before the replication package goes out.
\end{sowhat}
\end{frame}




\begin{frame}{A skill: \texttt{/explain-diff} }

\vspace{0.3cm}
 \texttt{/explain-diff} checks what changed within a Github repo and generates a html site that you can open

\begin{itemize}
  \vspace{0.2cm}
  \item \emphteal{Verification.} It diffs generated tables, compares $N$ and headline coefficients across logs, and quotes numbers only from output files it has opened.
  \vspace{0.2cm}
  \item \emphwarn{Output.} One offline HTML page in five sections: what the code did before, what changed and why, consequences for the results, a code walkthrough grouped by purpose, and a five-question quiz. Saved outside the repo with a plain-text summary in the chat.
\end{itemize}

\end{frame}

\begin{frame}{Dropbox or Github? }
\centering 
\begin{tabular}{>{\leavevmode\bfseries\color{mckdeep}}p{2.7cm}p{4.75cm}p{4.75cm}}
\toprule
\rowcolor{mcklight}
 & \textbf{GitHub} & \textbf{Dropbox} \\
\midrule
Sync      & Manual push and pull & Automatic, within seconds \\
History   & Full git history, branches & Dropbox revisions only \\
Best for  & Solo work, versioned drafts & Live coauthoring in Overleaf \\
Conflicts & Git merge & ``Conflicted copy'' files \\
\bottomrule
\end{tabular}

\vskip0.4cm
\begin{sowhat}
\emphwarn{Pick one per project.} Both syncs at once means duplicate-file chaos.
\end{sowhat}
\end{frame}

\begin{frame}{Two routes: GitHub Desktop for a paper trail (click to commit, no terminal), Dropbox for live coauthoring.}
\begin{columns}[t]
\column{0.49\textwidth}
\begin{mckcardg}[mcknavy]{Setup A\ \ GitHub via GitHub Desktop}{ovlB}
\scriptsize
\begin{enumerate}\setlength{\itemsep}{0pt}
  \item Overleaf: Menu $\to$ Integrations $\to$ GitHub $\to$ \emphnavy{Export Project to GitHub}
  \item GitHub Desktop: File $\to$ Clone repository, pick the repo and a local folder
  \item Open that folder in VS Code
  \item Edit in VS Code. In GitHub Desktop: one-line summary, \emphnavy{Commit to main}, \emphnavy{Push origin}
  \item Overleaf: Menu $\to$ GitHub $\to$ \emphnavy{Pull GitHub changes into Overleaf}
  \item Coauthor edited in Overleaf? \emphnavy{Push Overleaf changes to GitHub}, then pull in GitHub Desktop
\end{enumerate}
\textit{Accept the TeX \texttt{.gitignore} template: no \texttt{.pdf} or \texttt{.aux} in the history.}
\end{mckcardg}

\column{0.49\textwidth}
\begin{mckcardg}[mckteal]{Setup B\ \ Dropbox}{ovlB}
\scriptsize
\begin{enumerate}\setlength{\itemsep}{0pt}
  \item Install Dropbox and let the initial sync finish
  \item Overleaf: Menu $\to$ Dropbox Sync $\to$ \emphteal{Link}. Overleaf creates \texttt{Apps/Overleaf/} with one subfolder per project
  \item Open the project folder in VS Code and edit as normal --- every save reaches Overleaf within seconds, and Overleaf edits flow back
  \item Compile in Overleaf as usual, or locally with \texttt{latexmk -pdf}
\end{enumerate}
\textit{No commit step. Good for active coauthoring, less good if you want a paper trail.}
\end{mckcardg}
\end{columns}
\end{frame}


% ============================================================================
\section{Slides}
% ============================================================================
\sourceline{}

\begin{frame}{Making slides with AI}
\vskip0.2cm
\begin{columns}[t]
\column{0.48\textwidth}
\begin{mckcard}[mckgrey]{The old path}
\footnotesize
\begin{itemize}
  \item Open Beamer
  \item Write LaTeX by hand
  \item Fight the syntax
  \item Iterate slowly
\end{itemize}
\medskip
\textit{\textcolor{mckgrey}{Hours of mechanical work between idea and first draft.}}
\end{mckcard}

\column{0.48\textwidth}
\begin{mckcard}[mckblue]{The Claude Code path}
\footnotesize
\begin{itemize}
  \item Open the project folder
  \item Hand Claude a file: \textit{``Make a Beamer deck based on @paper.tex''}
  \item Compile, look at the output, iterate
\end{itemize}
\medskip
\textit{\textcolor{mckblue}{This deck was built that way.}}
\end{mckcard}
\end{columns}

\vskip0.4cm
\begin{sowhat}
This also means that we have to practice presentations in a different way
\end{sowhat}
\end{frame}

\begin{frame}{Including a style file gives Claude guidance to match the look you want.}
\vskip0.2cm
\textbf{Without a style file.} Claude defaults to generic Beamer themes. 

\vskip0.3cm
\textbf{With a style file.} The deck inherits your colour palette, frame structure, and title conventions.

\vskip0.4cm
\textbf{How to build a \texttt{slides-style.md}}
\begin{itemize}
  \item Ask Claude to analyze your old preentations or to make up something that looks good
  \item Save the result as \texttt{slides-style.md} in the project folder
  \item Reference it whenever you ask for a new deck
\end{itemize}

\vskip0.4cm
\begin{sowhat}
The Day 1 deck, this deck, and your future decks can all share one style file. Build it once.
\end{sowhat}
\end{frame}

% ============================================================================
\section{A replication package}
% ============================================================================
\sourceline{}

\begin{frame}{A replication package is a README, a master script, and portable paths.}
\begin{columns}[t]
\column{0.50\textwidth}
\textbf{What has to be in the folder}
\begin{enumerate}\small\setlength{\itemsep}{1pt}
  \item A README: data sources, access terms, software versions, runtime
  \item A table mapping each exhibit to the script that makes it and the file it writes
  \item One master script that runs the pipeline in order, relative paths, seeds set
  \item The data you may ship, and access instructions for the data you may not
\end{enumerate}

\small

\column{0.46\textwidth}
\begin{mckcardg}[mcknavy]{What to ask the agent}{repl}
\footnotesize
\begin{itemize}\setlength{\itemsep}{1pt}
  \item ``Read \texttt{Code/} and draft the README, including the exhibit-to-script table''
  \item ``List every input file the code reads; flag the ones not in this folder''
  \item ``Make every path relative and write \texttt{master.do}''
\end{itemize}
\vskip0.05em
\textcolor{mcknavy}{\itshape\scriptsize \textbf{Then verify:} run it in a fresh folder and diff the tables against the paper.}
\end{mckcardg}
\end{columns}

\begin{sowhat}
With register data the code and the README are the package --- the data never leaves the server.
\end{sowhat}
\end{frame}

\sourceline{}



% \begin{frame}{Live demo: Administrative tasks Claude Code can take off your desk.}
% \vskip0.2cm
% \textbf{Setup}
% \begin{itemize}\small
%   \item A folder with a call for papers
% %  \item Several CSVs of (fake) paper reviews
%  % \item Receipts, speaker lists, venue notes
% \end{itemize}

% \vskip0.3cm
% \textbf{The goal}
% \begin{itemize}\small
%   \item Build a budget for the conference
%  %\item Organise the reviews into rankings and themes
%  \end{itemize}

%  \vskip0.3cm
% \textbf{The output}
% \begin{itemize}\small
%   \item A budget spreadsheet with line items, totals, and assumptions
%  %\item Organise the reviews into rankings and themes
% \end{itemize}

% \vskip0.3cm
% \textbf{What we are using.} Plain Claude Code prompts --- no skills, no MCP. Just the folder and the agent.
% \end{frame}

% \begin{frame}{Two points the admin demo is meant to show }
% \vskip0.2cm
% \begin{enumerate}
%   \item \emphnavy{AI tools reduce the cost of administrative work.}
%   \begin{itemize}\small
%     \item Not just for research
%     \item Often the highest-leverage use case, because the friction is the highest and the quality bar is the lowest
%   \end{itemize}
%   \vskip0.3cm
%   \item \emphnavy{Context is the difference between useful and useless.}
%   \begin{itemize}\small
%     \item A chat tool gives you a generic conference-budget template
%     \item Claude Code finds your keynote speakers, your dinner, and your PhD travel funding by reading the folder
%   \end{itemize}
% \end{enumerate}

% \vskip0.4cm
% \begin{sowhat}
% The same pattern --- open the folder, hand the agent the materials, ask for the output --- works for grant reports, course design, referee assignments, and committee memos.
% \end{sowhat}
% \end{frame}


% \section{Literature}

% \begin{frame}
%   Convert files, 
% \end{frame}

% \begin{frame}{Zotero MCP connector}
  
% \end{frame}

% \begin{frame}{MCP servers plug the agent into external systems --- know it exists, install only with a clear use case.}
% \vskip0.3cm
% \textbf{Model Context Protocol.} A standard for letting Claude call external tools beyond the local filesystem.

% \vskip0.3cm
% \textbf{Examples already in production}
% \begin{itemize}
%   \item GitHub --- open PRs, read issues, comment on code review
%   \item Databases --- query Postgres, BigQuery, etc.\ directly
%   \item Custom servers you wrote yourself
% \end{itemize}

% \vskip0.3cm
% \textit{Like browser extensions, but for the agent. The list is growing fast.}

% \vskip0.3cm
% \begin{sowhat}
% For most academics, today, \emphnavy{you do not need MCP}. Know the concept exists; install one only when you have a concrete use case.
% \end{sowhat}
% \end{frame}

% 


%============================================================================
\section{Teaching in times of AI, continued}
% ============================================================================
\sourceline{Source: Bryan (2026), kevinbryanecon.com/teachwithai.html. }

\begin{frame}
\centering 
  \includegraphics[width=0.5\linewidth]{midterm_final_grades_crop}
\end{frame}

\begin{frame}{AI breaks the grade--knowledge link if we are not careful}
\vskip0.1cm
\textbf{Three things AI changes}
\begin{enumerate}\small\setlength{\itemsep}{2pt}
  \item \emphnavy{The grade--knowledge link is broken.} \\ A good take-home grade no longer shows the student can do the thing.
  \medskip
  \item \emphnavy{Students need to use AI well later.} \\
  Banning it everywhere leaves them untrained.
     \medskip
  \item \emphnavy{AI should let us teach more.} \\ 
  Tutoring, feedback, and individual practice used to be too expensive.
\end{enumerate}
\end{frame}


\begin{frame}{The pre-AI status quo was already leaky }
  US college study time fell from \textbf{40 to 27 hours a week}, 1961--2003
\begin{itemize}\setlength{\itemsep}{1pt}
  \item A quarter of finance students copied \textbf{wrong} Chegg answers, and learned less
\end{itemize}
\bigskip 
\textbf{And AI makes homework a worse signal}
\begin{itemize}\setlength{\itemsep}{1pt}
  \item Post-AI: homework scores $+$18\%, completion time $-$30\%, exam performance $-$20\% (China)
\end{itemize}

\end{frame}

\begin{frame}{How to protect learning? }
\begin{mckcardg}[mckteal]{Protect learning}{tchB}
\begin{enumerate}\setlength{\itemsep}{2pt}
  \item \textbf{Decide what should be learned without AI.} Be explicit about which assignments test AI-independent competence.
  \item \textbf{Give students an incentive to learn.} Regular, graded checks. Effort follows tests.
  \item \textbf{Do not give AI-cheatable assignments.} Every take-home worth real credit can be done by a frontier model today.
\end{enumerate}
\end{mckcardg}

\end{frame}

\begin{frame}{Using AI to teach more}

\begin{mckcardg}[mckteal]{Teach more with AI}{tchB}
\begin{enumerate}\setlength{\itemsep}{2pt}\setcounter{enumi}{3}
  \item \textbf{Help students learn more efficiently.} Class-specific tutoring, spaced repetition, mastery learning at scale.
  \item \textbf{Use AI to improve your own teaching.} Find out where you were confusing before the exam tells you.
  \item \textbf{Personalise assignments.} Problems at the level each student is struggling with.
  \item \textbf{Raise standards.} If research and drafting are cheaper, the bar moves up.
\end{enumerate}
\end{mckcardg}

\end{frame}




% \begin{frame}{Rules 1 and 4: audit what each assignment certifies and whether AI can do it.}
% \vskip0.1cm
% \begin{columns}[t]
% \column{0.46\textwidth}
% \textbf{How to run the audit}
% \begin{enumerate}\small\setlength{\itemsep}{1pt}
%   \item Syllabus, all assignments, one past exam in a folder
%   \item Ask Claude Code for the audit (right)
%   \item Make it \emph{actually attempt} one assignment. Grade it with your rubric
%   \item Decide per component: keep, move in-class, redesign
% \end{enumerate}

% \vskip0.1cm
% \textbf{Two questions per component}
% \begin{itemize}\footnotesize\setlength{\itemsep}{1pt}
%   \item What is it a \emph{proxy} for? Nobody inverts a matrix by hand, but being able to shows you understand the regression
%   \item Is it needed \emph{without} AI? A French essay written with AI is useless once translation is perfect
% \end{itemize}

% \column{0.52\textwidth}
% \begin{mckcardg}[mcknavy]{Prompt}{tchC}
% \footnotesize\itshape\color{mcknavy}
% ``Read the syllabus and every assignment in this folder. For each graded component, give me a table with: (1) the competency it is meant to certify, (2) whether students need that competency without AI, (3) whether a student using a frontier model could get a top mark on it today, and (4) what you would change. Be blunt.

% Then complete problem set 2 yourself, as a student would, and tell me what grade it would get.''
% \end{mckcardg}
% \footnotesize\textit{Expect the agent to score well. That is the point of the exercise.}
% \end{columns}

% \vskip0.2cm
% \begin{sowhat}
% Take-home work should no longer count, Bryan wrote on the day ChatGPT launched.
% \end{sowhat}
% % \end{frame}

% \begin{frame}{Three ways to make an assignment un-cheatable, and what each one costs you.}
% \vskip0.1cm
% \begin{columns}[t]
% \column{0.32\textwidth}
% \begin{mckcardg}[mcknavy]{1\ \ In-class only}{tchD}
% \footnotesize
% Oral exams (\emph{viva voce}), blue-book written exams, in-class problem solving.
% \medskip

% \textbf{Works because} AI is not in the room.
% \medskip

% \textcolor{mckwarn}{\textbf{Cost:} very time-intensive to evaluate. Hard to scale without TAs.}
% \end{mckcardg}

% \column{0.32\textwidth}
% \begin{mckcardg}[mckteal]{2\ \ Use AI to fight AI}{tchD}
% \footnotesize
% Students hand in a pr\'ecis early. You show the class what a frontier model produces from it. The final paper must be \emph{better than that}.
% \medskip

% \textbf{Works because} the AI output becomes the floor. To beat it: archives, phone calls, site visits, own data.
% \medskip

% \textcolor{mckwarn}{\textbf{Cost:} you generate and read the baseline. Term papers only.}
% \end{mckcardg}

% \column{0.32\textwidth}
% \begin{mckcardg}[mckamber]{3\ \ Un-cheatable at home}{tchD}
% \footnotesize
% Low-stakes quizzes graded 0 or 1. Pass with 10 or more correct, \emph{however many are wrong}. Wrong answers require the student to explain their logic.
% \medskip

% \textbf{Works because} more right answers buy nothing, so cheating buys nothing. Explaining errors is where the learning happens.
% \medskip

% \textcolor{mckwarn}{\textbf{Cost:} question bank and automation up front. Scales to any size.}
% \end{mckcardg}
% \end{columns}

% \vskip0.25cm
% \begin{sowhat}
% Bryan runs option 3 via his product, All Day TA. The design works with any quiz tool.
% \end{sowhat}
% \end{frame}

% \begin{frame}{The pr\'ecis workflow, step by step: make the model's paper the grade floor.}
% \vskip0.2cm
% \begin{tikzpicture}[
%   wstep/.style={rectangle, draw=mcknavy, line width=0.6pt, fill=mcknavy!7,
%     rounded corners=2pt, text width=3.1cm, minimum height=1.6cm, align=center,
%     font=\sffamily\footnotesize, text=mckdeep},
%   arr/.style={-{Stealth[length=2.5mm]}, line width=0.9pt, mckblue},
% ]
% \node[wstep] (s1) at (0,0)     {\textbf{1.\ Pr\'ecis, week 3}\\\scriptsize One page: question, argument, planned sources};
% \node[wstep] (s2) at (3.6,0)   {\textbf{2.\ Generate the baseline}\\\scriptsize Folder of pr\'ecis $\to$ agent writes the best paper a model can produce from each};
% \node[wstep] (s3) at (7.2,0)   {\textbf{3.\ Show the class}\\\scriptsize ``This is what AI gets you. Your grade starts here.''};
% \node[wstep] (s4) at (10.8,0)  {\textbf{4.\ Grade against it}\\\scriptsize Rubric line: value above the baseline. Primary sources, interviews, own data};
% \draw[arr] (s1) -- (s2);
% \draw[arr] (s2) -- (s3);
% \draw[arr] (s3) -- (s4);
% \end{tikzpicture}

% \vskip0.3cm
% \begin{columns}[t]
% \column{0.55\textwidth}
% \begin{mckcardg}[mcknavy]{Prompt for step 2}{tchE}
% \footnotesize\itshape\color{mcknavy}
% ``For each pr\'ecis in \texttt{precis/}, write the strongest 2,000-word term paper a frontier model can produce from it, using only what it already knows. Save as \texttt{baseline/<student>.md}. Then add a one-paragraph note on what a student would need to do to clearly beat this paper.''
% \end{mckcardg}

% \column{0.43\textwidth}
% \textbf{Practical notes}
% \begin{itemize}\footnotesize
%   \item Tell students exactly \emph{how} to beat the baseline. Bryan's list: get into the archives, get on the phone, visit the site, create data
%   \item Detection is a backstop, not the design. Pangram's false-positive rate is very low (Day 1), but the workflow works even if you never run it
%   \item Require an attached note on what AI did (see the policy slide)
% \end{itemize}
% \end{columns}
% \end{frame}

\begin{frame}{Students study when tested. Regular graded checks beat study plans.}

\begin{mckcardg}[mckteal]{How to implement it}{tchF}
  \begin{enumerate}\setlength{\itemsep}{2pt}
  \item Weekly quiz, pass/fail, small weight each, most weeks count
  \item Generate the bank from your own notes: \par
  {\itshape\color{mcknavy}``From \texttt{notes/week3.md}, write 15 multiple-choice questions at three levels, with distractors built from common misconceptions and a one-sentence explanation each. ''}
  \item (Verify every answer key yourself)
\end{enumerate}
\end{mckcardg}

\end{frame}

% \begin{frame}{Rules 5--7 rest on trial evidence: AI tutoring, teacher feedback, adaptive practice.}
% \vskip0.1cm
% \footnotesize
% \setlength{\tabcolsep}{4pt}
% \renewcommand{\arraystretch}{1.1}
% \begin{tabular}{>{\leavevmode\bfseries\color{mckdeep}}p{3.6cm}p{5.6cm}>{\leavevmode\color{mcknavy}\bfseries}p{3.5cm}}
% \toprule
% \rowcolor{mcklight}
% \textbf{\textcolor{mckdeep}{Idea}} & \textbf{\textcolor{mckdeep}{Evidence Bryan cites}} & \textbf{\textcolor{mckdeep}{Effect}} \\
% \midrule
% AI tutor (rule 5)
%   & Six-week trial with Nigerian students
%   & 0.31 SD \\
%   & World Bank meta-analysis of older ed-tech
%   & 0.12 SD \\
%   & Khanmigo (NBER); Harvard physics AI tutor
%   & Gains despite limited use \\
% Teacher feedback (rule 6)
%   & Feedback RCT, 1,000+ online CS teachers
%   & $+$13\% use of student ideas \\
% Adaptive practice (rule 7)
%   & Mindspark adaptive system (pre-LLM), 4 months
%   & 0.37 SD maths, 0.23 Hindi \\
%   & LLM-chosen problem sets (Taiwan), no-AI final
%   & 0.15 SD vs fixed sequence \\
% \midrule
% \textcolor{mckwarn}{Fetching answers from AI}
%   & Homework up, exams down (previous slides)
%   & \textcolor{mckwarn}{Does not work} \\
% \bottomrule
% \end{tabular}

% \vskip0.25cm
% \begin{sowhat}
% Old mechanisms, new cost: tutoring, mastery learning, spaced practice at scale.
% \end{sowhat}
% \end{frame}

\begin{frame}{Implementing tutors, misconception digests, levelled problem sets.}
\vskip0.1cm
\begin{columns}[t]
\column{0.32\textwidth}
\begin{mckcardg}[mcknavy]{Course tutor}{tchG}
\footnotesize
A shared assistant (custom GPT, Claude project, NotebookLM) loaded with \emph{your} lecture notes and problem sets.
\medskip

\textbf{System prompt essentials:} use the course's notation, never give the answer, ask the student to explain first, return to topics they missed.

\end{mckcardg}

\column{0.32\textwidth}
\begin{mckcardg}[mckteal]{Misconception digest}{tchG}
\footnotesize
Each week, export the questions students asked (tutor logs, forum, email), anonymised.
\medskip

\textbf{Prompt:} \textit{``Group these by underlying misconception, rank by frequency, and point to the lecture slide most likely to have caused each.''}

\end{mckcardg}

\column{0.32\textwidth}
\begin{mckcardg}[mckamber]{Levelled problem sets}{tchG}
\footnotesize
One problem, three difficulty tiers. Assign the tier by last week's quiz result.
\medskip

\textbf{Prompt:} \textit{``Write this problem at three levels: scaffolded, standard, extension. Same concept, full solutions, and a rubric I can grade against.''}
\medskip
\end{mckcardg}
\end{columns}

\vskip0.25cm
\begin{sowhat}
Student questions are personal data. Anonymise before exporting and check the institutional policy once.
\end{sowhat}
\end{frame}

% \begin{frame}{Rule 8: raise the bar, and write the AI policy into the syllabus.}
% \vskip0.1cm
% \begin{columns}[t]
% \column{0.55\textwidth}
% \begin{mckcardg}[mcknavy]{Bryan's course policy, verbatim}{tchH}
% \footnotesize\itshape
% ``You may use AI in this course. If you use AI to write, you must specify in an attached note precisely what AI assisted with. Understand that with AI, my expectation for the quality of your ideas and writing is higher than it used to be --- you need to show that you are able to provide value above and beyond AI. For readings, you must read them first yourself! Do not summarize them with AI. It will be very obvious in class if you do. I want your ideas; I can ask GPT for its thoughts without coming to class, as can your classmates.''
% \end{mckcardg}

% \column{0.43\textwidth}
% \textbf{Why raise the bar}
% \begin{itemize}\footnotesize\setlength{\itemsep}{1pt}
%   \item Randomised workplace studies: AI raised writing quality and cut completion time
%   \item Consultants working inside AI's competence produced better work \citep{dell2026navigating}
%   \item If content-gathering is cheap, an A from ten years ago is no longer an A
% \end{itemize}

% \vskip0.15cm
% \textbf{Practical}
% \begin{itemize}\footnotesize\setlength{\itemsep}{1pt}
%   \item State the policy \emph{per assignment}, not once per course
%   \item Require the attached AI-use note and grade it
%   \item Rubric line: sloppy citation and error-ridden prose now cost more, not less
% \end{itemize}
% \end{columns}

% \vskip0.2cm
% \begin{sowhat}
% The clearer you are about where students \emphnavy{should} use AI and where they \emphnavy{must not}, the better they get the balance right.
% \end{sowhat}
% \end{frame}







% ----------------------------------------------------------------------------
% CLOSING
%----------------------------------------------------------------------------
\sourceline{}

\section{Conclusions }

\begin{frame}
  \centering
  \textbf{\emphwarn{Key message:} Context improves AI tools tremendously, and the goal is to provide the right context} \\
  \bigskip 
  \textbf{\emphwarn{Key difference:} how you provide the useful context and how easy it is}
  \bigskip \end{frame}


\begin{frame}[plain]
\begin{tikzpicture}[remember picture, overlay]
  \fill[mckdeep] (current page.south west) rectangle (current page.north east);
  \node[anchor=center, text=white, font=\Huge\sffamily\bfseries]
    at (current page.center)
    {Q\,\&\,A};
  \node[anchor=center, text=mckblue!70!white, font=\normalsize\sffamily\itshape,
        text width=0.7\paperwidth, align=center]
    at ([yshift=-1.2cm]current page.center)
    {Open floor. Email: \texttt{claes.backman@gmail.com}\ \textbar\
     Slides, skills, guides: \texttt{claesbackman.com/lugano}};
  \draw[mckblue, line width=2pt]
    ([xshift=-1.25cm, yshift=0.6cm]current page.center)
    -- ([xshift=1.25cm, yshift=0.6cm]current page.center);
\end{tikzpicture}
\end{frame}

% \begin{frame}[allowframebreaks]{References}
% \footnotesize
% \bibliographystyle{plainnat}
% \bibliography{references}
% \end{frame}

% ============================================================================
\appendix
\sourceline{}



% ============================================================================
\section{Setting up your workspace}
% ============================================================================
\sourceline{}

\begin{frame}{Five steps get VS Code + Claude Code working in under five minutes.}
\vskip0.2cm
\begin{enumerate}
  \item Open VS Code and the Extensions panel (\texttt{Ctrl+Shift+X} / \texttt{Cmd+Shift+X})
  \item Search \emphnavy{``Claude Code''} and install \emph{Claude Code for VS Code} (extension ID \texttt{anthropic.claude-code})
  \item Restart VS Code --- Claude Code opens in the sidebar
  \item Sign in with your Claude account (Pro/Max/Team/Enterprise all include Claude Code) or an API key
  \item Open your \emphnavy{full project folder} (File $>$ Open Folder) --- never a single file
\end{enumerate}

\vskip0.3cm
\textbf{What an IDE actually is.} Editor + file explorer + terminal + extensions, in one window. VS Code is the dominant free one. Think of it as Word or R Studio for code, LaTeX, markdown, and plain text.

% \vskip0.2cm
% \begin{sowhat}
% The file tree on the left shows you what Claude is reading; the diff view shows you exactly what it changed. These two windows are why VS Code beats the terminal for academic work.
% \end{sowhat}
\end{frame}


\begin{frame}{How to get started}

  Get the skills from \href{https://claesbackman.com/lugano}{\texttt{claesbackman.com/lugano}} --- download the zip, or paste the one-line terminal command on that page. Source and updates: \href{https://github.com/claesbackman/AI-research-feedback}{github.com/claesbackman/AI-research-feedback}. \\
  \bigskip 

  Unzip, then copy each skill folder (\texttt{review-paper}, \texttt{review-paper-light}, \ldots) into \texttt{\textasciitilde/.claude/skills/} --- a hidden folder under your user folder (Windows: \texttt{C:\textbackslash Users\textbackslash you\textbackslash .claude\textbackslash skills})
  \begin{itemize}
    \item To show hidden folders: Open File Explorer, click the View tab at the top, select Show, and check Hidden items
    \item A skill in \texttt{\textasciitilde/.claude/skills/} works in every project; one in \texttt{project/.claude/skills/} only in that project
  \end{itemize}
  \bigskip
  Restart VS Code\\
    \bigskip
  The skill loads when you type \texttt{/skill-name}
\end{frame}

\begin{frame}{A handful of extensions cover Stata, Python, LaTeX, and data}
\vskip0.2cm
\small
\begin{tabular}{>{\bfseries\color{mckdeep}}l>{\ttfamily}lp{6.0cm}}
\toprule
\rowcolor{mcklight}
\textbf{\textcolor{mckdeep}{Stack}} & \normalfont\textbf{\textcolor{mckdeep}{Extension ID}} & \textbf{\textcolor{mckdeep}{What it does}} \\
\midrule
Stata  & tmonk.stata-workbench    & Run Stata from the editor  \\
Python & ms-python.python         & Syntax, IntelliSense, run button \\
LaTeX  & james-yu.latex-workshop  & Compile on save, side-by-side PDF preview, inline errors \\
Data   & mechatroner.rainbow-csv  & Colour-coded CSV columns; SQL-like filtering \\
       & tomoki1207.pdf            & Renders PDFs inline --- read codebooks next to your code \\
\bottomrule
\end{tabular}

\vskip0.3cm
\begin{sowhat}
Install Python even if you do not write it. The agent uses Python as glue for many different tasks. 
\end{sowhat}
\end{frame}

\sourceline{Source: claesbackman.com/overleaf-vscode-guide.html; Overleaf documentation.}


\begin{frame}{You can link Overleaf to your own computer and work in that folder}

\textbf{Why bother?} Overleaf has no Claude Code, no Codex, no Copilot, no proper diff view, no terminal. The fix: write in VS Code and let Overleaf follow what you produce.

\medskip
\textbf{Requirement.} Both routes need \emphnavy{Overleaf Premium} --- many universities provide it through the institutional login.
\end{frame}



% \begin{frame}{Six things that will catch you out when syncing Overleaf.}
% \footnotesize
% \renewcommand{\arraystretch}{1.0}
% \begin{tabular}{>{\leavevmode\bfseries\color{mckwarn}}p{3.0cm}p{4.9cm}p{4.6cm}}
% \toprule
% \rowcolor{mcklight}
% \textbf{\textcolor{mckdeep}{Gotcha}} & \textbf{\textcolor{mckdeep}{What happens}} & \textbf{\textcolor{mckdeep}{What to do}} \\
% \midrule
% The first sync overwrites
%   & If one side is empty, the other side wins --- not always in the direction you expect
%   & Make sure Overleaf and the repo or folder agree on contents before linking \\
% GitHub sync is not instant
%   & Overleaf only pulls when someone clicks
%   & Tell coauthors to hit \emphnavy{Pull GitHub changes} after you push \\
% Dropbox conflicts are silent
%   & Two people edit the same file and Dropbox creates \texttt{paper (conflicted copy).tex} without warning
%   & Skim the project folder occasionally \\
% Dropbox notifications are noisy
%   & Auto-sync produces a constant stream of update messages
%   & Mute them --- or use GitHub (this is why I switched) \\
% Local compile needs packages
%   & Compiles on Overleaf, fails locally: Overleaf ships a large TeX Live
%   & \texttt{tlmgr install \textless name\textgreater} for whatever is missing \\
% Figure paths must stay inside the project
%   & Overleaf rejects parent-directory references like \texttt{../Figures/fig1.pdf}
%   & Keep figures in a subfolder of the project root \\
% \bottomrule
% \end{tabular}
%
% \vskip0.2cm
% \textit{Once the paper is local, the agents edit the \texttt{.tex} directly --- that is the whole point of the detour.}
% \end{frame}

% ============================================================================
% \section{Working with the agent}
% ============================================================================
\sourceline{}

\begin{frame}{Every Claude Code session: open the folder, set context, then prompt.}
\vskip0.2cm
\begin{enumerate}
  \item Open VS Code and the project folder (File $\to$ Open Folder)
  \item Open the Claude Code panel in the sidebar --- or the terminal (\texttt{Ctrl+\char`\`}) and type \texttt{claude}
  \item Confirm Claude is pointed at the \emphnavy{full project folder} --- not a single file
  \item Send a first orienting prompt
\end{enumerate}

\vskip0.4cm
\textbf{First prompt to try:}
\begin{quote}
\itshape\color{mcknavy}
``Read the files in this folder and tell me what this project is about.''
\end{quote}

\vskip0.2cm
\textbf{Expected behaviour.} Claude reads the folder, summarises the structure, and asks you what you want to do next. 
\end{frame}


\begin{frame}[fragile]{\texttt{CLAUDE.md} is the standing-instructions file that is always loaded}
\vskip0.1cm
\small
\begin{verbatim}
# About me
I am a researcher in [field]. I write papers, lecture notes, policy pieces.

# How I want you to work with me
- Ask clarifying questions before generating long output.
- Keep a skeptical, critical tone in feedback. Do not flatter.
- When editing, preserve my voice. Avoid generic LLM phrasing.
- Avoid bullet lists and passive voice in formal writing.
- Cite the specific line or section you are commenting on.

# Things to avoid
- Do not fabricate citations.
- Do not over-claim causality when summarising empirical work.
- Do not insert emoji or markdown decorations in formal documents.
\end{verbatim}

\vskip0.2cm
% \begin{sowhat}
% Loaded for the project, every session. The single highest-leverage one-hour investment you can make.
% \end{sowhat}
\end{frame}

\begin{frame}{A few files hold your durable project context}
\vskip0.1cm
\footnotesize
\begin{tabular}{>{\leavevmode\bfseries\color{mckdeep}}p{3.9cm}p{8.5cm}}
\toprule
\rowcolor{mcklight}
\textbf{\textcolor{mckdeep}{File}} & \textbf{\textcolor{mckdeep}{Use it for}} \\
\midrule
\texttt{project/CLAUDE.md} & Project instructions. Dataset structure, sample restrictions, paper title, coauthor conventions, which script produces which output. \\
\addlinespace
\texttt{\textasciitilde/.claude/CLAUDE.md} & Personal preferences across every project: writing style, default tone, citation conventions. Project instructions are loaded later and usually carry more weight.\\
\addlinespace
\texttt{voice.md}         & Writing samples. Paragraphs, abstracts, and tone notes that Claude should match when editing prose. \\
\addlinespace
\texttt{.claude/skills/...} & Reusable workflows. Named tasks you want to run again without rewriting the prompt. \\
\bottomrule
\end{tabular}

\vskip0.15cm
\begin{sowhat}
Keep secrets and raw confidential data \emphwarn{out of all of them} --- these files are read into context on every session and may be sent to the model provider.
\end{sowhat}
\end{frame}

\begin{frame}{Bootstrap \texttt{CLAUDE.md} in two minutes --- type \texttt{/init}.}
\vskip0.2cm
\textbf{The built-in command:} type \texttt{/init} in the Claude Code panel. Or spell it out:
\begin{quote}
\itshape\color{mcknavy}
``Read this project folder and draft a \texttt{CLAUDE.md} file with the project structure, conventions, and instructions you should follow. Do not edit anything else yet.''
\end{quote}

\vskip0.3cm
\textbf{What Claude does}
\begin{itemize}
  \item Reads the folder
  \item Infers languages, structure, conventions
  \item Generates a \texttt{CLAUDE.md} draft
\end{itemize}

\vskip0.3cm
\textbf{Then you edit it.} Add tone preferences, things to avoid, sample restrictions Claude could not infer.

% \vskip0.3cm
% \begin{sowhat}
% Two minutes to a working \texttt{CLAUDE.md}. The draft is the starting point, not the final word.
% \end{sowhat}
\end{frame}

% \begin{frame}{Plan mode --- think before you write. The agent plans, you correct, then it executes.}
% \vskip0.2cm
% For anything non-trivial, ask Claude to outline its approach \emphnavy{before} generating.

% \vskip0.3cm
% \textbf{Example prompt:}
% \begin{quote}
% \itshape\color{mcknavy}
% ``Before writing the policy brief, outline the structure, the three findings you plan to lead with, and the one caveat you'll flag. Do not write the brief yet. Ask me questions if you need to.''
% \end{quote}

% \vskip0.3cm
% You review the plan, correct the framing, then approve it. Claude Code has an explicit \emphnavy{plan mode}:
% \begin{itemize}
%   \item Accept the plan and it executes; reject it and you have lost thirty seconds
% \end{itemize}

% \vskip0.3cm
% \begin{sowhat}
% Planning catches the wrong approach when it is cheap to fix --- before the agent has written 800 lines you have to read.
% \end{sowhat}
% \end{frame}

\begin{frame}{Permissions follow a tiered model}
\vskip0.3cm
\begin{tikzpicture}[
  tier/.style={rounded corners=4pt, draw, line width=0.8pt,
    minimum width=13cm, minimum height=1.0cm,
    inner sep=0pt, anchor=west},
  tname/.style={anchor=west, font=\sffamily\bfseries\small, text width=3.0cm},
  tdesc/.style={anchor=west, font=\sffamily\small, text=mckgrey, text width=8.5cm},
]
\node[tier, fill=mckteal!10, draw=mckteal] (t1) at (0,0) {};
\node[tname, text=mckteal!60!black] at ([xshift=0.3cm]t1.west) {Plan mode};
\node[tdesc] at ([xshift=3.5cm]t1.west) {Reads and proposes; cannot write.};

\node[tier, fill=mcknavy!10, draw=mcknavy,
      below=4pt of t1.south west, anchor=north west] (t2) {};
\node[tname, text=mcknavy] at ([xshift=0.3cm]t2.west) {Manual};
\node[tdesc] at ([xshift=3.5cm]t2.west) {Asks before each edit or command.};

\node[tier, fill=mckamber!10, draw=mckamber,
      below=4pt of t2.south west, anchor=north west] (t3) {};
\node[tname, text=mckamber] at ([xshift=0.3cm]t3.west) {Edit automatically};
\node[tdesc] at ([xshift=3.5cm]t3.west) {Claude will edit your selected text or the whole file};

\node[tier, fill=mckwarn!10, draw=mckwarn,
      below=4pt of t3.south west, anchor=north west] (t4) {};
\node[tname, text=mckwarn] at ([xshift=0.3cm]t4.west) {Auto mode (default mode now)};
\node[tdesc] at ([xshift=3.5cm]t4.west) {Claude will approve actions that pass a safety check and pause for anything risky};
\end{tikzpicture}

\vskip0.5cm
Claude Code should stay inside the folder you opened, but this is not certain. Use \texttt{/permissions} to see the current rules, and allow- or block-list specific commands (\texttt{Rscript}, \texttt{pdflatex}, \texttt{rm}) in \texttt{settings.json}. 
\end{frame}

\begin{frame}{A session is a conversation, not a query}
\vskip0.1cm
\begin{center}
\begin{tikzpicture}[
  box/.style={rectangle, draw=mcknavy, line width=0.6pt, fill=mcknavy!7, rounded corners=2pt,
    text width=3.0cm, minimum height=1.0cm, align=center, font=\sffamily\footnotesize, text=mckdeep},
  kill/.style={rectangle, draw=mckwarn, line width=0.6pt, fill=mckwarn!8, rounded corners=2pt,
    text width=3.6cm, minimum height=1.0cm, align=center, font=\sffamily\footnotesize, text=mckdeep},
  arr/.style={-{Stealth[length=2.5mm]}, line width=0.9pt, mckblue},
  lbl/.style={font=\sffamily\scriptsize, text=mckgrey},
]
\node[box] (p) at (0,0)   {\textbf{1.\ Prompt}\\\scriptsize specific, with the context it needs};
\node[box] (r) at (4.6,0) {\textbf{2.\ Response}\\\scriptsize a draft, not an answer};
\node[box] (j) at (9.2,0) {\textbf{3.\ You react}\\\scriptsize accept, correct, redirect};
\draw[arr] (p) -- (r);
\draw[arr] (r) -- (j);
\draw[arr] (j.north) -- ++(0,0.45) -| (p.north) node[lbl, pos=0.25, above] {next turn builds on your corrections --- 10--20 turns is normal};
\node[kill] (k) at (9.2,-1.75) {\textbf{Off track?}\\\scriptsize kill it and restart with a better first prompt};
\draw[arr, mckwarn] (j.south) -- (k.north);
\node[lbl, anchor=east, text width=5.2cm, align=right] at ([xshift=-0.4cm]k.west) {Signs: repeats a mistake, argues in circles, forgets an earlier instruction, three turns spent fixing one thing};
\end{tikzpicture}
\end{center}

\vskip0.3cm
\footnotesize
\textbf{Continue the conversation.} The context carries your corrections, so later answers improve: \textit{``Shorter, and use the Danish case.''} \textit{``You misread Table 2 --- look again.''} \textit{``Now the same for Section 3.''}

\medskip
\textbf{Kill it when you veer off track.} Start fresh and put what you learned into the first prompt --- or into \texttt{CLAUDE.md} so you never repeat it. 

\end{frame}


\begin{frame}[fragile]{Appendix --- \texttt{CLAUDE.md} starter template.}
\small
\begin{verbatim}
# About me
I am a researcher in [field].

# How I want you to work with me
- Ask clarifying questions before generating long output.
- Critical, skeptical tone in feedback. Do not flatter.
- When editing, preserve my voice (see voice.md). No generic LLM phrasing.
- Avoid bullet lists and passive voice in formal writing.
- Cite the specific line or section you are commenting on.

# Things to avoid
- Do not fabricate citations.
- Do not over-claim causality.
- Do not insert emoji or markdown decorations in formal docs.
\end{verbatim}
\end{frame}

\begin{frame}{Appendix --- instructor's skill library.}
\vskip0.2cm
\small
\begin{tabular}{>{\bfseries\color{mckdeep}}ll}
\toprule
\rowcolor{mcklight}
\textbf{\textcolor{mckdeep}{Skill}} & \textbf{\textcolor{mckdeep}{Use}} \\
\midrule
\texttt{/review-paper}       & Pre-submission referee report, 8 agents \\
\texttt{/review-paper-light} & Fast 2-agent version, $\approx$ 1 minute \\
\texttt{/review-paper-code}  & Paper + code alignment review \\
\texttt{/review-grant}       & Pre-submission grant review, 6 agents \\
\texttt{/review-pap}         & Pre-analysis-plan review \\
\texttt{/audit-analysis}     & Adversarial audit of changed analysis code \\
\texttt{/explain-diff}       & Explain a code change as an HTML page with a quiz \\
\texttt{/paper-version}      & Paper $\to$ policy brief / 1-page / 5-page \\
\texttt{/pdf-to-markdown}    & Convert PDFs to readable markdown \\
\bottomrule
\end{tabular}

\vskip0.3cm
\textit{All at \href{https://github.com/claesbackman/AI-research-feedback}{github.com/claesbackman/AI-research-feedback} (MIT), linked from \href{https://claesbackman.com/lugano}{\texttt{claesbackman.com/lugano}}. Fork them, modify them, make them yours.}
\end{frame}

\begin{frame}{Appendix --- file format quick reference.}
\vskip0.2cm
\small
\begin{tabular}{>{\bfseries\color{mckdeep}}lp{8.1cm}}
\toprule
\rowcolor{mcklight}
\textbf{\textcolor{mckdeep}{Format}} & \textbf{\textcolor{mckdeep}{Handling}} \\
\midrule
\texttt{.md}, \texttt{.tex}, \texttt{.txt}, \texttt{.bib} & Native --- no conversion \\
\texttt{.py}, \texttt{.R}, \texttt{.do}, \texttt{.ipynb}  & Native --- read, edit, run \\
\texttt{.csv}                                              & Native up to a few thousand rows \\
\texttt{.docx}                                             & Convert: \texttt{pandoc file.docx -o file.md} \\
\texttt{.xlsx}                                             & Export to CSV first \\
\texttt{.dta}, \texttt{.rds}                               & Script-based conversion to CSV \\
PDF (born-digital, short)                                  & Native; expect table mangling \\
PDF (long, scanned, equations)                             & Convert to markdown first \\
\bottomrule
\end{tabular}
\end{frame}

\begin{frame}{Appendix --- signs of session drift.}
\vskip0.3cm
\begin{itemize}
  \item Claude ignores \texttt{CLAUDE.md}
  \item Repeats mistakes you already corrected
  \item Invents files that do not exist
  \item Gives circular answers
  \item Forgets the last three turns
\end{itemize}

\vskip0.4cm
\textbf{Fix.} Start a fresh session or close and reopen. Re-state the task. Re-attach the relevant file.

\vskip0.3cm
\textit{30+ turns or the task has shifted $\to$ start fresh. Always.}
\end{frame}

\end{document}
